For a type \(\tau=(\delta,\eta)\), define its slack at \(\lambda^\dagger\) by
\[s_{\delta,\eta}=\sum_{z=1}^{3}\lambda^\dagger_{(\eta(\delta(z)),\delta(z),z)}-\bigl(\eta(1)-\eta(0)\bigr).\]
Using the within-arm coordinate order \(((0,0),(1,0),(0,1),(1,1))\), direct substitution gives the complete table
\[\begin{array}{c|rrrr}\delta\backslash\eta&00&01&10&11\\ \hline 000&2&1&1&0\\001&2&2&0&0\\010&0&1&1&2\\011&0&2&0&2\\100&2&0&2&0\\101&2&1&1&0\\110&0&0&2&2\\111&0&1&1&2\end{array}.\]
Every entry is nonnegative, and the fixed coordinates in arms two and three are zero. Thus \(\lambda^\dagger\in\mathcal H^0_{2,3}\). Reading off every zero gives exactly the signature asserted in the statement.

Stack the ten coefficient rows indexed, in the displayed order, by \(B^\dagger\), and append the two gauge rows. In the ambient order obtained by concatenating the three arm blocks, the resulting matrix is
\[M^\dagger=\left(\begin{array}{rrrrrrrrrrrr}0&1&0&0&0&1&0&0&0&1&0&0\\0&1&0&0&0&1&0&0&0&0&1&0\\0&1&0&0&0&1&0&0&0&0&0&1\\1&0&0&0&0&0&1&0&1&0&0&0\\1&0&0&0&0&0&1&0&0&0&1&0\\0&1&0&0&0&0&1&0&0&0&1&0\\0&0&0&1&1&0&0&0&1&0&0&0\\0&0&0&1&0&1&0&0&0&1&0&0\\0&0&1&0&0&0&1&0&1&0&0&0\\0&0&0&1&0&0&0&1&1&0&0&0\\0&0&0&0&1&0&0&0&0&0&0&0\\0&0&0&0&0&0&0&0&1&0&0&0\end{array}\right).\]
Integer elimination gives \(\det(M^\dagger)=1\). Hence the first ten rows have rank \(10\), while all twelve have rank \(12\). The gauge section has dimension \(10\), so the ten tight equations isolate \(\lambda^\dagger\) in that section. If \(\lambda^\dagger\) were a nontrivial midpoint of two feasible points, every inequality tight at the midpoint would also be tight at both endpoints; the endpoints would then belong to the same isolated affine solution set and would coincide. Therefore \(\lambda^\dagger\) is a vertex.

We next verify the literal comparator family. By the gauge-transversal lemma, adding an element of \(K_{2,3}\) adds a constant within each arm, and therefore
\[I_z(\lambda+v)=\lambda_{(1,1,z)}+a_z-\lambda_{(0,1,z)}-a_z=I_z(\lambda).\]
At \(n=2\), an admissible index \((a,s)\) has \(s=(s_0,s_1)\), where \(a,s_0,s_1\) are the three distinct arms. Directly from the treatment-swapped formula,
\[I_z(\bar w(a,s))=-w(a,s)_{1,0,z}+w(a,s)_{0,0,z}=\mathbf1\{z=s_0\}.\]
The correction made by \(\Gamma\) is armwise constant, so \(I(\Gamma(\bar w(a,s)))=e_{s_0}\).

For \(a=1\) and \(s=(2,3)\), the gauge-fixed representative has blocks
\[\mu^0=\bigl((1,0,1,1),(0,0,-1,0),(0,0,0,0)\bigr).\]
Its complete slack table is
\[\begin{array}{c|rrrr}\delta\backslash\eta&00&01&10&11\\ \hline 000&1&0&1&0\\001&1&0&1&0\\010&0&0&0&0\\011&0&0&0&0\\100&1&0&2&1\\101&1&0&2&1\\110&0&0&1&1\\111&0&0&1&1\end{array}.\]
Thus \(\mu^0\) is feasible. In the reduced coordinate order
\[x=(u_{0,1},u_{1,1},v_{0,1},v_{1,1},u_{1,2},v_{0,2},v_{1,2},u_{1,3},v_{0,3},v_{1,3}),\]
the tight types \(000/01,000/11,001/01,001/11,010/00,010/01,010/10,011/00,100/01,110/00\) have row-normal matrix
\[R_0=\begin{pmatrix}1&0&0&0&0&0&0&0&0&0\\0&1&0&0&1&0&0&1&0&0\\1&0&0&0&0&0&0&0&0&1\\0&1&0&0&1&0&0&0&0&1\\1&0&0&0&0&1&0&0&0&0\\1&0&0&0&0&0&1&0&0&0\\0&1&0&0&0&1&0&1&0&0\\1&0&0&0&0&1&0&0&1&0\\0&0&0&1&0&0&0&0&0&0\\0&0&1&0&0&1&0&0&0&0\end{pmatrix},\qquad\det(R_0)=-1.\]
The same active-rank argument makes \(\mu^0\) a vertex. Every admissible ordered triple \((a,s_0,s_1)\) is an arm permutation of \((1,2,3)\); the arm-permutation lemma therefore makes all six comparator representatives vertices. Define the additional gauge-invariant contrast
\[J_z(\lambda)=\lambda_{(1,0,z)}-\lambda_{(0,0,z)}.\]
The literal formula gives
\[J(\Gamma(\bar w(a,s)))=-e_a,\qquad I(\Gamma(\bar w(a,s)))=e_{s_0}.\]
Thus the pair \((J,I)\) recovers \((a,s_0)\), and \(s_1\) is the remaining arm. The six indices give six distinct comparator vertices.

Finally, direct evaluation at the displayed witness gives
\[I(\lambda^\dagger)=(-1,2,1),\qquad J(\lambda^\dagger)=(-1,-2,1).\]
For the three arms, the ordered pairs \((J_z,I_z)\) are respectively
\[(-1,-1),\quad(-2,2),\quad(1,1),\]
and are pairwise distinct. The arm-permutation lemma shows that \(T_\sigma\lambda^\dagger\) is a vertex for every \(\sigma\in S_3\), and that its ordered triple of \((J,I)\)-pairs is the corresponding permutation of these three distinct pairs. No nonidentity permutation fixes that ordered triple, so the orbit has all six elements. Each of the three original arm blocks is nonconstant, and arm permutation and gauge correction preserve this property; hence every orbit member has three essential blocks. Its \(I\)-vector is a coordinate permutation of \((-1,2,1)\), which cannot be a standard basis vector. The six orbit vertices are therefore disjoint from the six comparator vertices. It follows that the complement has at least six elements and that the whole vertex set has at least \(6+6=12\) elements, proving the remaining conclusions.